\cotefile{93}

\scanpage{193}

que lorsque il n'est pas nul, il \uline{n'est pas identique} au groupe des $\omega$ dont un multiple est logarithmique, i.e. il y a des \deleted{\unsure{des str.}} connexions intégrables sur $\mathcal O_X$ qui ne sont pas isomorphes à la connexion standard, [en fait le groupe des classes d'isomorphie de telles connexions intégrables est isomorphe à \deleted{$\mathbb Q/\mathbb Z$} $L\otimes\mathbb Q/\mathbb Z\simeq(\mathbb Q/\mathbb Z)^{r}$]. Par contre, si $X$ est complet, \inserted{donc $\Gamma(X,\mathcal O_X^{*})$ \deleted{\unsure{formé seul.}} formé de sections horizontales,} \unsure{tte|une} connexion intégrable sur $\mathcal O_X$ est triviale, i.e. une $1$-forme différentielle qui est loc. ét une dérivée logarithmique, est nulle. \deleted{La classification des modules inversibles avec \uline{connexion à courbure} nulle, i.e. à stratification, est donc}

Revenons au cas où $k$ est de car. quelconque. Je dis que si $\Gamma(X,\mathcal O_X)^{*}$ est formé de fonctions horizontales, \deleted{\unsure{des}} p.ex. $X$ propre sur $k$ et séparable sur $k$, alors sur $\underline O_X$ il n'y a comme seule \deleted{\unsure{structure}} stratification \uline{intégrable} que la structure habituelle. En effet, $\underline O_X^{\otimes n}=\underline O_X$ \inserted{($n$ premier à $p$)} a une section horizontale \inserted{$\sigma$, pr $\unsure{n\ne0|n\gg0|\text{tout}\ne0}$}, \deleted{\unsure{donc qui est}} pour la stratification $T^{\otimes n}$, \deleted{qui est aussi horizontale}. Donc la section $\sigma_0=1$ de $\underline O_X$ satisfait $\sigma_0^{\otimes n}=\lambda\sigma$, $\lambda\in\Gamma(X,\underline O_X)^{*}$, \deleted{donc} et $\lambda$ horizontale, \inserted{par hyp.,} donc $\lambda\sigma$ aussi, donc $\sigma_0^{\otimes n}$ horizontale, donc $\sigma_0$ aussi. // Par suite, on voit que la classification des \deleted{\unsure{fibr.}} modules inversibles à connexion intégrable est, dans ce cas, \uline{identique} à celle des modules inversibles \deleted{\unsure{dont}} $L$ tels que $\exists\ n>0$, $n$ premier à $p$, avec $L^{\otimes n}\simeq\underline O_X$

% archive sheet number: 192
